\documentclass[11pt,a4paper]{article}
\usepackage{fontspec,amsmath,amssymb,amsthm,microtype}
\setmainfont{lmroman10-regular.otf}[BoldFont=lmroman10-bold.otf,ItalicFont=lmroman10-italic.otf,BoldItalicFont=lmroman10-bolditalic.otf]
\usepackage{xcolor}
\usepackage[most]{tcolorbox}
\usepackage{mathrsfs}
\providecommand{\lt}{<}
\providecommand{\gt}{>}
\setmonofont{DejaVuSansMono.ttf}[BoldFont=DejaVuSansMono-Bold.ttf,ItalicFont=DejaVuSansMono-Oblique.ttf,BoldItalicFont=DejaVuSansMono-BoldOblique.ttf,Scale=0.86]
\newfontfamily\leanfallback{STIXTwoMath-Regular.otf}[Scale=0.86]
\definecolor{leanrule}{gray}{0.62}
\definecolor{leanback}{gray}{0.962}
\definecolor{leantitle}{gray}{0.22}
\definecolor{leanmuted}{gray}{0.38}
\newtcolorbox{leanbox}[1]{enhanced,breakable,colback=leanback,colframe=leanrule,
  boxrule=0.4pt,arc=0.6mm,left=2.4mm,right=2.4mm,top=1.8mm,bottom=1.8mm,
  before skip=0.8em,after skip=1em,fonttitle=\small\bfseries,coltitle=leantitle,
  colbacktitle=leanback,titlerule=0.3pt,title={#1}}
\newlength{\leancharwidth}
\newenvironment{leancode}{\par\smallskip\begingroup\ttfamily\footnotesize\raggedright
  \settowidth{\leancharwidth}{M}\setlength{\parskip}{0pt}\setlength{\parindent}{0pt}
  \frenchspacing}{\par\endgroup\smallskip}
\newcommand{\leanline}[2]{\hangindent=\dimexpr#1\leancharwidth+4\leancharwidth\relax
  \hangafter=1\noindent\hspace*{#1\leancharwidth}#2\par}
\newcommand{\leansource}[1]{\par{\footnotesize\color{leanmuted}#1}\par}
\newcommand{\leanlabel}[1]{\par\smallskip\noindent{\small\bfseries #1}\par}
\usepackage[margin=25mm]{geometry}
\usepackage[unicode,hidelinks]{hyperref}
\hypersetup{pdfauthor={LangevinLean Project}}
\setlength{\emergencystretch}{3em}
\setlength{\parskip}{0.4em}
\setlength{\parindent}{0pt}
\allowdisplaybreaks[2]
\raggedbottom
\title{\texorpdfstring{Data processing and a finite-entropy change of reference}{Data processing and a finite-entropy change of reference}}
\author{LangevinLean Project}
\date{}
\begin{document}\maketitle
Draft mathematical exposition; these notes may contain mistakes.


Apply the same Gaussian transition to the initial law \(\mu\) and the Gibbs law \(\nu\). Relative-entropy data processing gives \(H(\mu P_t\mid\nu P_t)\le H(\mu\mid\nu)\). To compare \(\mu P_t\) with \(\nu\), use the density bound for \(\nu P_t\) relative to \(\nu\).



The logarithmic likelihood through an intermediate reference is the sum of the two logarithmic likelihoods wherever their positive finite density ratios are defined. An integrable upper bound controls its positive part, and the exponential of the negative likelihood controls its negative part. This proves finite entropy against the new reference before the integrated comparison is used. The Gaussian reference-density bound is a constant plus a multiple of \(\|\nabla U\|^2\), and propagation of the gradient moment makes that correction integrable without an initial position moment.


\section{\texorpdfstring{KL is unchanged by adjoining the same independent law}{KL is unchanged by adjoining the same independent law}}


Adjoin the same independent probability innovation \(\eta\) to two finite measures. The log likelihood of \(\mu\otimes\eta\) relative to \(\nu\otimes\eta\) depends only on the first coordinate, so KL is unchanged. The proof uses the standard conditional-product KL identity with the constant Markov kernel \(x\mapsto\eta\); its conditional products are exactly the ordinary product measures. This equality supplies the independent-noise step in data processing.



Taking a product with the same independent probability does not change the likelihood ratio: it depends only on the first coordinate. Integrating out the extra factor therefore preserves the relative entropy.

\begin{leanbox}{Lean: \texttt{klDiv\_\allowbreak{}prod\_\allowbreak{}same\_\allowbreak{}right}}
\leansource{\texttt{klDiv\_\allowbreak{}prod\_\allowbreak{}same\_\allowbreak{}right} \quad UnadjustedLangevin/\allowbreak{}LSI/\allowbreak{}EntropyEndpoint.lean, lines 25–29}
\begin{leancode}
\leanline{0}{theorem\ klDiv\_prod\_same\_right\ (μ\ ν\ :\ Measure\ Ω)\ [IsFiniteMeasure\ μ]\ [IsFiniteMeasure\ ν]}
\leanline{4}{(η\ :\ Measure\ Ξ)\ [IsProbabilityMeasure\ η]\ :}
\leanline{4}{InformationTheory.klDiv\ (μ.prod\ η)\ (ν.prod\ η)\ =\ InformationTheory.klDiv\ μ\ ν}
\end{leancode}
\leanlabel{How the Lean proof follows the argument}
\texttt{kl\allowbreak{}Div\_\allowbreak{}prod\_\allowbreak{}same\_\allowbreak{}right}: the proof is the application \texttt{simpa only [Measure.\allowbreak{}comp\allowbreak{}Prod\_\allowbreak{}const] using}. It specializes that result with the hypotheses and objects in the exact statement.
\end{leanbox}

\section{\texorpdfstring{Adding log likelihoods through an intermediate law}{Adding log likelihoods through an intermediate law}}


For \(\mu\ll\rho\ll\nu\), the Radon-Nikodym chain rule gives \(d\mu/d\nu=(d\mu/d\rho)(d\rho/d\nu)\) almost everywhere under the integrated law \(\mu\). Each factor is positive and finite there: positivity is transported along absolute continuity, and finiteness uses the probability measures' finite Radon-Nikodym derivatives.



On this full-measure set the density ratios are ordinary positive finite values, so the logarithm product rule applies. This proves \(\log(d\mu/d\nu)=\log(d\mu/d\rho)+\log(d\rho/d\nu)\), as an almost-everywhere identity under \(\mu\). The positivity and finiteness checks prevent using logarithm algebra at a zero or infinite density.



For three mutually compatible absolutely continuous laws, their density ratios multiply. Taking logarithms gives the sum of the two log likelihoods on the common full-measure set, before any integration is attempted.

\begin{leanbox}{Lean: \texttt{llr\_\allowbreak{}eq\_\allowbreak{}add\_\allowbreak{}of\_\allowbreak{}absolutelyContinuous}}
\leansource{\texttt{llr\_\allowbreak{}eq\_\allowbreak{}add\_\allowbreak{}of\_\allowbreak{}absolutelyContinuous} \quad UnadjustedLangevin/\allowbreak{}LSI/\allowbreak{}EntropyEndpoint.lean, lines 32–45}
\begin{leancode}
\leanline{0}{theorem\ llr\_eq\_add\_of\_absolutelyContinuous}
\leanline{4}{(μ\ ρ\ ν\ :\ Measure\ Ω)\ [IsProbabilityMeasure\ μ]\ [IsProbabilityMeasure\ ρ]}
\leanline{4}{[IsProbabilityMeasure\ ν]\ (hμρ\ :\ μ\ {\leanfallback ≪}\ ρ)\ (hρν\ :\ ρ\ {\leanfallback ≪}\ ν)\ :}
\leanline{4}{llr\ μ\ ν\ =ᵐ[μ]\ fun\ x\ ↦\ llr\ μ\ ρ\ x\ +\ llr\ ρ\ ν\ x}
\end{leancode}
\leanlabel{How the Lean proof follows the argument}
\texttt{llr\_\allowbreak{}eq\_\allowbreak{}add\_\allowbreak{}of\_\allowbreak{}absolutely\allowbreak{}Continuous}: the intermediate fact is \texttt{h\ensuremath{\mu}\ensuremath{\nu}}, in source order; the source steps apply or rewrite \texttt{Pi.\allowbreak{}mul\_\allowbreak{}apply}, \texttt{ENNReal.\allowbreak{}to\allowbreak{}Real\_\allowbreak{}mul}, \texttt{ht.\allowbreak{}ne}, \texttt{hq.\allowbreak{}ne'}, \texttt{hu.\allowbreak{}ne}.
\end{leanbox}

\section{\texorpdfstring{An exponential bound on the negative log likelihood}{An exponential bound on the negative log likelihood}}


For \(\mu\ll\nu\), the reciprocal likelihood \(e^{-\log(d\mu/d\nu)}\) agrees \(\mu\)-almost everywhere with the real reverse Radon-Nikodym derivative \((d\nu/d\mu)\). This derivative is integrable under \(\mu\), even when \(\nu\) has a singular component relative to \(\mu\). Thus the negative exponential of the log likelihood is always integrable; finite KL has not been assumed at this stage.



The exponential of the negative log likelihood is a reciprocal density ratio. Its integral under the first law is bounded by the reference mass, giving control of the negative log-likelihood tail.

\begin{leanbox}{Lean: \texttt{integrable\_\allowbreak{}exp\_\allowbreak{}neg\_\allowbreak{}llr}}
\leansource{\texttt{integrable\_\allowbreak{}exp\_\allowbreak{}neg\_\allowbreak{}llr} \quad UnadjustedLangevin/\allowbreak{}LSI/\allowbreak{}EntropyEndpoint.lean, lines 49–52}
\begin{leancode}
\leanline{0}{theorem\ integrable\_exp\_neg\_llr}
\leanline{4}{(μ\ ν\ :\ Measure\ Ω)\ [IsProbabilityMeasure\ μ]\ [IsProbabilityMeasure\ ν]}
\leanline{4}{(hμν\ :\ μ\ {\leanfallback ≪}\ ν)\ :\ Integrable\ (fun\ x\ ↦\ Real.exp\ (-llr\ μ\ ν\ x))\ μ}
\end{leancode}
\leanlabel{How the Lean proof follows the argument}
\texttt{integrable\_\allowbreak{}exp\_\allowbreak{}neg\_\allowbreak{}llr}: the source steps apply or rewrite \texttt{Measure.\allowbreak{}integrable\_\allowbreak{}to\allowbreak{}Real\_\allowbreak{}rn\allowbreak{}Deriv.\allowbreak{}congr}, \texttt{exp\_\allowbreak{}neg\_\allowbreak{}llr}.
\end{leanbox}

\section{\texorpdfstring{Log-likelihood integrability from an upper bound}{Log-likelihood integrability from an upper bound}}


Suppose the log likelihood \(\ell=\log(d\mu/d\nu)\) has an integrable upper bound \(b\). Where \(\ell\ge0\), \(|\ell|\le b\le|b|\). Where \(\ell<0\), the elementary exponential bound gives \(|\ell|=-\ell\le e^{-\ell}\). Hence globally \(|\ell|\le|b|+e^{-\ell}\). The first term is integrable by hypothesis, and the second by reciprocal-likelihood integrability. Measurable domination proves \(\ell\in L^1(\mu)\), controlling its negative tail without a lower density bound.



An integrable upper bound controls the positive part, while the exponential estimate controls the negative part. Together they establish absolute integrability of the log likelihood needed for the entropy chain rule.

\begin{leanbox}{Lean: \texttt{integrable\_\allowbreak{}llr\_\allowbreak{}of\_\allowbreak{}integrable\_\allowbreak{}upper\_\allowbreak{}bound}}
\leansource{\texttt{integrable\_\allowbreak{}llr\_\allowbreak{}of\_\allowbreak{}integrable\_\allowbreak{}upper\_\allowbreak{}bound} \quad UnadjustedLangevin/\allowbreak{}LSI/\allowbreak{}EntropyEndpoint.lean, lines 56–70}
\begin{leancode}
\leanline{0}{theorem\ integrable\_llr\_of\_integrable\_upper\_bound}
\leanline{4}{(μ\ ν\ :\ Measure\ Ω)\ [IsProbabilityMeasure\ μ]\ [IsProbabilityMeasure\ ν]}
\leanline{4}{(hμν\ :\ μ\ {\leanfallback ≪}\ ν)\ (b\ :\ Ω\ →\ ℝ)\ (hb\ :\ Integrable\ b\ μ)}
\leanline{4}{(hbound\ :\ ∀ᵐ\ x\ ∂μ,\ llr\ μ\ ν\ x\ ≤\ b\ x)\ :\ Integrable\ (llr\ μ\ ν)\ μ}
\end{leancode}
\leanlabel{How the Lean proof follows the argument}
\texttt{integrable\_\allowbreak{}llr\_\allowbreak{}of\_\allowbreak{}integrable\_\allowbreak{}upper\_\allowbreak{}bound}: the cases are separated by \texttt{by\_\allowbreak{}cases hp : 0 \ensuremath{\le} llr \ensuremath{\mu} \ensuremath{\nu} x}; the intermediate facts include \texttt{he}, \texttt{hn}, in source order; the source steps apply or rewrite \texttt{hb.\allowbreak{}norm.\allowbreak{}add}, \texttt{integrable\_\allowbreak{}exp\_\allowbreak{}neg\_\allowbreak{}llr}, \texttt{Real.\allowbreak{}exp\_\allowbreak{}nonneg}, \texttt{abs\_\allowbreak{}of\_\allowbreak{}neg}, \texttt{lt\_\allowbreak{}of\_\allowbreak{}not\_\allowbreak{}ge}. The remaining scalar rearrangements are checked by \texttt{linarith}.
\end{leanbox}

\section{\texorpdfstring{Entropy comparison through a reference density bound}{Entropy comparison through a reference density bound}}


Assume finite \(H(\mu\mid\rho)\), \(\rho\ll\nu\), and \(d\rho/d\nu\le e^b\), with \(b\in L^1(\mu)\). Finite KL gives \(\mu\ll\rho\), hence \(\mu\ll\nu\). The density bound transfers to \(\mu\)-almost every point, where positivity of the intermediate density permits logarithms. Thus \(\log(d\rho/d\nu)\le b\).



The log-likelihood chain rule gives \(\log(d\mu/d\nu)\le\log(d\mu/d\rho)+b\). Its right side is integrable. The preceding one-sided integrability theorem therefore proves the new log likelihood integrable and hence \(H(\mu\mid\nu)<\infty\). Only after proving finiteness do we integrate the inequality, yielding \(H(\mu\mid\nu)\le H(\mu\mid\rho)+\mathbb E_\mu b\). Finite entropy against the new reference is a conclusion, rather than an additional premise.



Decompose the log likelihood through the intermediate reference law. Its first entropy is controlled by data processing, and the density-ratio upper bound controls the added term by an integrable observable.

\begin{leanbox}{Lean: \texttt{klDiv\_\allowbreak{}finite\_\allowbreak{}and\_\allowbreak{}le\_\allowbreak{}of\_\allowbreak{}reference\_\allowbreak{}density\_\allowbreak{}bound}}
\leansource{\texttt{klDiv\_\allowbreak{}finite\_\allowbreak{}and\_\allowbreak{}le\_\allowbreak{}of\_\allowbreak{}reference\_\allowbreak{}density\_\allowbreak{}bound} \quad UnadjustedLangevin/\allowbreak{}LSI/\allowbreak{}EntropyEndpoint.lean, lines 75–103}
\begin{leancode}
\leanline{0}{theorem\ klDiv\_finite\_and\_le\_of\_reference\_density\_bound}
\leanline{4}{(μ\ ρ\ ν\ :\ Measure\ Ω)\ [IsProbabilityMeasure\ μ]\ [IsProbabilityMeasure\ ρ]}
\leanline{4}{[IsProbabilityMeasure\ ν]}
\leanline{4}{(hKL\ :\ InformationTheory.klDiv\ μ\ ρ\ ≠\ ⊤)\ (hρν\ :\ ρ\ {\leanfallback ≪}\ ν)}
\leanline{4}{(b\ :\ Ω\ →\ ℝ)\ (hb\ :\ Integrable\ b\ μ)}
\leanline{4}{(hbound\ :\ ∀ᵐ\ x\ ∂ν,\ (ρ.rnDeriv\ ν\ x).toReal\ ≤\ Real.exp\ (b\ x))\ :}
\leanline{4}{InformationTheory.klDiv\ μ\ ν\ ≠\ ⊤\ ∧}
\leanline{6}{(InformationTheory.klDiv\ μ\ ν).toReal\ ≤}
\leanline{8}{(InformationTheory.klDiv\ μ\ ρ).toReal\ +\ ∫\ x,\ b\ x\ ∂μ}
\end{leancode}
\leanlabel{How the Lean proof follows the argument}
\texttt{kl\allowbreak{}Div\_\allowbreak{}finite\_\allowbreak{}and\_\allowbreak{}le\_\allowbreak{}of\_\allowbreak{}reference\_\allowbreak{}density\_\allowbreak{}bound}: the intermediate facts include \texttt{h\ensuremath{\mu}\ensuremath{\nu}}, \texttt{hchain}, \texttt{hupper}, \texttt{hnew}, \texttt{hi}, in source order; the source steps apply or rewrite \texttt{Real.\allowbreak{}log\_\allowbreak{}le\_\allowbreak{}iff\_\allowbreak{}le\_\allowbreak{}exp}, \texttt{ENNReal.\allowbreak{}to\allowbreak{}Real\_\allowbreak{}pos}, \texttt{Information\allowbreak{}Theory.\allowbreak{}to\allowbreak{}Real\_\allowbreak{}kl\allowbreak{}Div\_\allowbreak{}of\_\allowbreak{}measure\_\allowbreak{}eq}, \texttt{Pi.\allowbreak{}add\_\allowbreak{}apply}, \texttt{integral\_\allowbreak{}add}. The remaining scalar rearrangements are checked by \texttt{linarith}.
\end{leanbox}

\section{\texorpdfstring{Data processing for the Gaussian transition}{Data processing for the Gaussian transition}}


The interpolation is the measurable image of the initial law times standard Gaussian noise under \((x,z)\mapsto x-t\nabla U(x)+\sqrt{2t}z\). Apply KL data processing to this same map for both initial laws. The independent-noise product identity removes the common Gaussian factor, giving \(H(\rho_t^\mu\mid\rho_t^\nu)\le H(\mu\mid\nu)\). This argument compares two evolved laws and does not assert that either law is invariant under the discrete step.



Both inputs are passed through the same measurable Gaussian update with the same independent noise law. Relative-entropy data processing therefore bounds the two output laws by the divergence of their inputs.

\begin{leanbox}{Lean: \texttt{interpolationLaw\_\allowbreak{}klDiv\_\allowbreak{}le}}
\leansource{\texttt{interpolationLaw\_\allowbreak{}klDiv\_\allowbreak{}le} \quad UnadjustedLangevin/\allowbreak{}LSI/\allowbreak{}EntropyEndpoint.lean, lines 110–119}
\begin{leancode}
\leanline{0}{theorem\ interpolationLaw\_klDiv\_le}
\leanline{4}{(U\ :\ E\ →\ ℝ)\ (L\ :\ ℝ≥0)\ (hLip\ :\ LipschitzWith\ L\ (gradient\ U))}
\leanline{4}{(μ\ ν\ :\ Measure\ E)\ [IsProbabilityMeasure\ μ]\ [IsProbabilityMeasure\ ν]\ (t\ :\ ℝ)\ :}
\leanline{4}{InformationTheory.klDiv\ (interpolationLaw\ U\ μ\ t)\ (interpolationLaw\ U\ ν\ t)\ ≤}
\leanline{6}{InformationTheory.klDiv\ μ\ ν}
\end{leancode}
\leanlabel{How the Lean proof follows the argument}
\texttt{interpolation\allowbreak{}Law\_\allowbreak{}kl\allowbreak{}Div\_\allowbreak{}le}: the intermediate fact is \texttt{h}, in source order; the source steps apply or rewrite \texttt{kl\allowbreak{}Div\_\allowbreak{}prod\_\allowbreak{}same\_\allowbreak{}right}.
\end{leanbox}


Earlier results used here: \cite[Section~3.36.5]{LangevinLean2026ULASC}; \cite[Section~3.36.7]{LangevinLean2026ULASC}.

\section{\texorpdfstring{Propagation of the gradient moment}{Propagation of the gradient moment}}


Write \(\nabla U(Y_t)=[\nabla U(Y_t)-\nabla U(X)]+\nabla U(X)\). The initial gradient is \(L^2\), and the frozen-error energy theorem makes the bracketed discrepancy \(L^2\) on the independent initial/\allowbreak{}noise space. Their sum is \(L^2\); the map-measure equivalence transfers this to the interpolation law.



Pointwise, the squared norm of the sum is at most twice each squared norm. Integrate and insert \(\mathbb E\|\nabla U(Y_t)-\nabla U(X)\|^2\le L^2(t^2G_0+2td)\). The initial-coordinate marginal gives the remaining \(G_0\), so \(G_t\le2(1+L^2t^2)G_0+4tdL^2\). This proves both integrability and the quantitative estimate without a position moment.



Lipschitzness bounds \(\|\nabla U(Y_t)\|^2\) by \(2\|\nabla U(X)\|^2+2L^2\|Y_t-X\|^2\). The exact displacement moment then yields the displayed propagated energy.

\begin{leanbox}{Lean: \texttt{interpolationLaw\_\allowbreak{}gradient\_\allowbreak{}moment}}
\leansource{\texttt{interpolationLaw\_\allowbreak{}gradient\_\allowbreak{}moment} \quad UnadjustedLangevin/\allowbreak{}LSI/\allowbreak{}EntropyEndpoint.lean, lines 123–182}
\begin{leancode}
\leanline{0}{theorem\ interpolationLaw\_gradient\_moment}
\leanline{4}{(U\ :\ E\ →\ ℝ)\ (L\ :\ ℝ≥0)\ (hLip\ :\ LipschitzWith\ L\ (gradient\ U))}
\leanline{4}{(μ\ :\ Measure\ E)\ [IsProbabilityMeasure\ μ]}
\leanline{4}{(hi\ :\ Integrable\ (fun\ x\ ↦\ ‖gradient\ U\ x‖\ \textasciicircum{}\ 2)\ μ)}
\leanline{4}{(t\ :\ ℝ)\ (ht\ :\ 0\ ≤\ t)\ :}
\leanline{4}{Integrable\ (fun\ x\ ↦\ ‖gradient\ U\ x‖\ \textasciicircum{}\ 2)\ (interpolationLaw\ U\ μ\ t)\ ∧}
\leanline{6}{(∫\ x,\ ‖gradient\ U\ x‖\ \textasciicircum{}\ 2\ ∂interpolationLaw\ U\ μ\ t)\ ≤}
\leanline{8}{2\ *\ (1\ +\ (L\ :\ ℝ)\ \textasciicircum{}\ 2\ *\ t\ \textasciicircum{}\ 2)\ *\ (∫\ x,\ ‖gradient\ U\ x‖\ \textasciicircum{}\ 2\ ∂μ)\ +}
\leanline{10}{4\ *\ t\ *\ (Module.finrank\ ℝ\ E\ :\ ℝ)\ *\ (L\ :\ ℝ)\ \textasciicircum{}\ 2}
\end{leancode}
\leanlabel{How the Lean proof follows the argument}
\texttt{interpolation\allowbreak{}Law\_\allowbreak{}gradient\_\allowbreak{}moment}: the intermediate facts include \texttt{hgrad}, \texttt{hg0}, \texttt{hproposal}, \texttt{hg0eq}, \texttt{houteq}, in source order; the source steps apply or rewrite \texttt{he.\allowbreak{}add}, \texttt{hgrad.\allowbreak{}comp\_\allowbreak{}fst}, \texttt{prob\allowbreak{}Real\_\allowbreak{}univ}, \texttt{one\_\allowbreak{}smul}, \texttt{integral\_\allowbreak{}map\_\allowbreak{}of\_\allowbreak{}strongly\allowbreak{}Measurable}. The remaining scalar rearrangements are checked by \texttt{nlinarith}.
\end{leanbox}


Earlier results used here: \cite[Section~3.34.1]{LangevinLean2026ULASC}; \cite[Section~3.36.5]{LangevinLean2026ULASC}; \cite[Section~3.36.7]{LangevinLean2026ULASC}.

\section{\texorpdfstring{The interpolation entropy upper bound}{The interpolation entropy upper bound}}


Let \(\rho_t\) be the evolved input and \(\sigma_t\) the evolved Gibbs target. Data processing gives \(H(\rho_t\mid\sigma_t)\le H_0<\infty\). Suppose the reference law satisfies \(d\sigma_t/d\nu\le\exp(b+c\|\nabla U\|^2)\), with its absolute continuity retained explicitly.



Finite initial KL first supplies finite \(G_0\) through the direct entropy/\allowbreak{}score estimate. The propagated gradient-moment theorem then proves \(G_t<\infty\). Consequently the logarithmic reference correction \(b+c\|\nabla U\|^2\) is integrable under \(\rho_t\). Apply the finite-entropy change-of-reference theorem to obtain \(H(\rho_t\mid\nu)<\infty\) and \(H(\rho_t\mid\nu)\le H_0+b+cG_t\). The preceding Gibbs comparison establishes the pointwise reference-density bound.



Insert the explicit log density-ratio bound and the propagated gradient moment into the entropy decomposition. This gives a finite upper bound on interpolation entropy, suitable for sending elapsed time to zero.

\begin{leanbox}{Lean: \texttt{interpolationLaw\_\allowbreak{}entropy\_\allowbreak{}bound\_\allowbreak{}of\_\allowbreak{}density\_\allowbreak{}bound}}
\leansource{\texttt{interpolationLaw\_\allowbreak{}entropy\_\allowbreak{}bound\_\allowbreak{}of\_\allowbreak{}density\_\allowbreak{}bound} \quad UnadjustedLangevin/\allowbreak{}LSI/\allowbreak{}EntropyEndpoint.lean, lines 187–222}
\begin{leancode}
\leanline{0}{theorem\ interpolationLaw\_entropy\_bound\_of\_density\_bound}
\leanline{4}{(U\ :\ E\ →\ ℝ)\ (L\ :\ ℝ≥0)\ (hL\ :\ 0\ <\ L)\ (hU\ :\ Differentiable\ ℝ\ U)}
\leanline{4}{(hLip\ :\ LipschitzWith\ L\ (gradient\ U))}
\leanline{4}{[IsProbabilityMeasure\ (ULA.gibbsMeasure\ U)]}
\leanline{4}{(μ\ :\ Measure\ E)\ [IsProbabilityMeasure\ μ]}
\leanline{4}{(hKL\ :\ InformationTheory.klDiv\ μ\ (ULA.gibbsMeasure\ U)\ ≠\ ⊤)}
\leanline{4}{(t\ :\ ℝ)\ (ht\ :\ 0\ ≤\ t)\ (b\ c\ :\ ℝ)}
\leanline{4}{(href\ :\ interpolationLaw\ U\ (ULA.gibbsMeasure\ U)\ t\ {\leanfallback ≪}\ ULA.gibbsMeasure\ U)}
\leanline{4}{(hbound\ :\ ∀ᵐ\ x\ ∂ULA.gibbsMeasure\ U,}
\leanline{6}{((interpolationLaw\ U\ (ULA.gibbsMeasure\ U)\ t).rnDeriv\ (ULA.gibbsMeasure\ U)\ x).toReal\ ≤}
\leanline{8}{Real.exp\ (b\ +\ c\ *\ ‖gradient\ U\ x‖\ \textasciicircum{}\ 2))\ :}
\leanline{4}{InformationTheory.klDiv\ (interpolationLaw\ U\ μ\ t)\ (ULA.gibbsMeasure\ U)\ ≠\ ⊤\ ∧}
\leanline{6}{(InformationTheory.klDiv\ (interpolationLaw\ U\ μ\ t)\ (ULA.gibbsMeasure\ U)).toReal\ ≤}
\leanline{8}{(InformationTheory.klDiv\ μ\ (ULA.gibbsMeasure\ U)).toReal\ +\ b\ +}
\leanline{10}{c\ *\ (∫\ x,\ ‖gradient\ U\ x‖\ \textasciicircum{}\ 2\ ∂interpolationLaw\ U\ μ\ t)}
\end{leancode}
\leanlabel{How the Lean proof follows the argument}
\texttt{interpolation\allowbreak{}Law\_\allowbreak{}entropy\_\allowbreak{}bound\_\allowbreak{}of\_\allowbreak{}density\_\allowbreak{}bound}: the intermediate facts include \texttt{hdpi}, \texttt{hbetween}, \texttt{hreal}, in source order; the source steps apply or rewrite \texttt{integral\_\allowbreak{}add}, \texttt{integrable\_\allowbreak{}const}, \texttt{integral\_\allowbreak{}const}, \texttt{prob\allowbreak{}Real\_\allowbreak{}univ}, \texttt{one\_\allowbreak{}smul}. The remaining scalar rearrangements are checked by \texttt{linarith}.
\end{leanbox}


Earlier results used here: \cite[Section~3.26.3]{LangevinLean2026ULASC}.

\begin{thebibliography}{ULA-LSI}
\bibitem[ULA-SC]{LangevinLean2026ULASC}
LangevinLean Project. \emph{Unadjusted Langevin under Strong Convexity}. LL-01, 2026. Draft report, content revision d50c0f64815b. \url{https://odakyildiz.github.io/LangevinLean/docs/pdf/UnadjustedLangevin/FullProof.pdf}.
\end{thebibliography}

\end{document}
